\documentclass[11pt,a4paper]{article}

% ── Packages ──────────────────────────────────────────────────────────────────
\usepackage[margin=1in]{geometry}
\usepackage{amsmath,amssymb}
\usepackage{graphicx}
\usepackage{booktabs}
\usepackage{hyperref}
\usepackage{xcolor}
\usepackage{enumitem}
\usepackage{float}
\usepackage{multirow}
\usepackage{array}
\usepackage{caption}
\usepackage{subcaption}
\usepackage{fancyhdr}
\usepackage{titlesec}
\usepackage{parskip}

\hypersetup{
    colorlinks=true,
    linkcolor=blue!70!black,
    urlcolor=blue!70!black,
    citecolor=blue!70!black
}

\pagestyle{fancy}
\fancyhf{}
\rhead{2026 Quantum Sustainability Challenge}
\lhead{Wildfire Risk \& Insurance Premium Modeling}
\rfoot{Page \thepage}

% ── Title ─────────────────────────────────────────────────────────────────────
\title{
    \vspace{-1cm}
    \LARGE\textbf{Wildfire Risk and Insurance Premium Modeling} \\[0.4em]
    \large 2026 Quantum Sustainability Challenge Submission
}
\author{
    \textbf{[YOUR NAME(S)]} \\
    Team: [TEAM NAME] \\
    [UNIVERSITY / AFFILIATION] \\
    \texttt{[EMAIL ADDRESS]}
}
\date{}

% ══════════════════════════════════════════════════════════════════════════════
\begin{document}
\maketitle
\thispagestyle{fancy}

% ── Abstract (max 400 words) ──────────────────────────────────────────────────
\begin{abstract}
\noindent
We present a hybrid Quantum Machine Learning (QML) framework for wildfire risk
modeling across California ZIP codes, addressing both Task~1 (fire risk prediction
for 2023) and Task~2 (insurance premium forecasting). Using historical wildfire
and climate data spanning 2018--2022 across 2,593 California ZIP codes,
we formulate wildfire prediction as a zero-inflated count regression problem:
the expected number of fires per ZIP code per month.

For Task~1A, we implement a Variational Quantum Circuit (VQC) with 9~qubits
using angle encoding and data re-uploading. Nine carefully engineered features
capture the primary drivers of California wildfire risk: temperature, precipitation,
dryness accumulation, seasonality, and geography. The model is trained to predict
$\log(1 + n_{\text{fires}})$ and outputs a continuous monthly fire risk score per ZIP
code, which is aggregated into an annual risk score by summing across months.

For Task~1B, we benchmark the QML approach against five classical models
(Dummy, Ridge, Random Forest, Gradient Boosting, and XGBoost) using a strict
temporal train/validate/test split (2018--2020 / 2021 / 2022). The best
classical model, Gradient Boosting, achieves an AUC-ROC of 0.873 on the
held-out 2022 test set and correctly ranks fire-prone ZIP-month pairs 87\%
of the time. The QML model is evaluated on the same split and metrics for
a direct comparison.

For Task~2, annual fire risk scores from our model are used as inputs to a
time series insurance premium model predicting 2021 premiums from 2018--2020
historical data.

Our data pipeline, feature engineering, and complete model code are
available at: \href{https://github.com/[YOUR-REPO]}{github.com/[YOUR-REPO]}.
\end{abstract}

\tableofcontents
\newpage

% ══════════════════════════════════════════════════════════════════════════════
\section{Data Description and Preprocessing}
\label{sec:data}

\subsection{Source Dataset}
The provided wildfire dataset contains monthly records for California ZIP codes
covering the period 2017--2023. Each record includes climate variables
(daily maximum and minimum temperature, total precipitation), a derived dryness
index, and binary fire occurrence flags with counts of wildfire events.
The dataset was provided as \texttt{cleaned\_data\_sorted.csv}.

\subsection{Data Cleaning}
The following preprocessing steps were applied, implemented in \texttt{clean\_data.py}:

\begin{enumerate}[leftmargin=*]
    \item \textbf{Year filtering.} Records were filtered to the training window
          2018--2022, yielding 125,164 rows before gap-filling.

    \item \textbf{Deduplication.} The raw data contained 444 duplicate
          ZIP$\times$year$\times$month entries arising from two sources:
          (i)~the same fire event geo-coded to multiple coordinate points within
          the same ZIP, and (ii)~separate rows for weather data and fire event
          data for the same period. These were resolved by keeping the
          weather-valid row, setting \texttt{has\_fire}~=~1 if any row in the
          group indicated a fire, and summing fire counts into \texttt{num\_fires}.

    \item \textbf{Complete panel construction.} A full panel of
          $2{,}593\ \text{ZIP codes} \times 5\ \text{years} \times 12\ \text{months}
          = 155{,}580$ rows was constructed. ZIP$\times$year$\times$month
          combinations absent from the raw data were inserted with
          \texttt{has\_fire}~=~0 and \texttt{num\_fires}~=~0 (no fire assumed).

    \item \textbf{Temperature imputation.} Missing values in
          \texttt{avg\_tmax\_c} (31,158 rows) and \texttt{avg\_tmin\_c} (31,158
          rows) were filled using the mean for that ZIP$\times$month combination
          across the available years. \texttt{temp\_mean} and
          \texttt{temp\_range} were recomputed from the filled values.

    \item \textbf{Precipitation imputation.} All 31,116 missing
          \texttt{tot\_prcp\_mm} values (entirely in 2022) were filled using
          the ZIP$\times$month historical mean across 2018--2021.

    \item \textbf{Dryness index imputation.} Missing \texttt{dryness\_index}
          values (30,860 rows) were recomputed using the formula
          \begin{equation}
              \text{dryness\_index} = \frac{T_{\max}}{P + 1}
              \label{eq:dryness}
          \end{equation}
          where $T_{\max}$ is \texttt{avg\_tmax\_c} and $P$ is
          \texttt{tot\_prcp\_mm}. The $+1$ smoothing term prevents division
          by zero and recovers $T_{\max}$ as the dryness measure when there
          is no precipitation, consistent with all observed values in the data.

    \item \textbf{Geocoding.} Latitude and longitude were populated for all
          2,593 ZIP codes using the \texttt{pgeocode} library, replacing the
          largely missing coordinate columns in the raw data.
\end{enumerate}

The cleaned dataset (\texttt{cleaned\_data\_final.csv}) contains 155,580 rows
across 14 columns with no missing values in any feature used for modeling.

\subsection{Class Distribution}
\label{sec:imbalance}

The dataset is severely zero-inflated. Of the 155,580 ZIP-month observations,
153,962 (99.0\%) have \texttt{num\_fires}~=~0. Only 1,618 observations (1.0\%)
record at least one fire. Among fire observations, the distribution of
\texttt{num\_fires} is:

\begin{table}[H]
\centering
\caption{Distribution of \texttt{num\_fires} across all ZIP-month observations.}
\begin{tabular}{rr}
\toprule
\textbf{num\_fires} & \textbf{Rows} \\
\midrule
0 & 153,962 \\
1 & 1,404 \\
2 & 165 \\
3 & 33 \\
4 & 9 \\
5 & 5 \\
6 & 2 \\
\midrule
\textbf{Total} & \textbf{155,580} \\
\bottomrule
\end{tabular}
\label{tab:fire_dist}
\end{table}

This imbalance is the central challenge of the modeling task and motivates both
the choice of AUC-ROC as the primary evaluation metric and the use of
sample-weighted training.

% ══════════════════════════════════════════════════════════════════════════════
\section{Feature Engineering}
\label{sec:features}

Nine features were selected and engineered to capture the primary physical
drivers of California wildfire risk while remaining within a circuit budget
of 9 qubits for the QML model. Redundant features were excluded:
\texttt{temp\_mean} was dropped because it is an exact linear combination of
\texttt{avg\_tmax\_c} and \texttt{avg\_tmin\_c}; \texttt{dryness\_index}
was dropped because it is fully determined by \texttt{avg\_tmax\_c} and
\texttt{tot\_prcp\_mm} via Equation~\eqref{eq:dryness}, and the QML
circuit can learn this interaction through entangling gates.

\subsection{Feature Definitions}

\begin{table}[H]
\centering
\caption{Final feature set and rationale.}
\begin{tabular}{lll}
\toprule
\textbf{Feature} & \textbf{Source} & \textbf{Rationale} \\
\midrule
\texttt{month\_sin} & $\sin(2\pi m/12)$ & \multirow{2}{*}{Cyclical month encoding — December and January are adjacent} \\
\texttt{month\_cos} & $\cos(2\pi m/12)$ & \\
\texttt{year} & raw & Linear climate trend (2018$\to$2022) \\
\texttt{avg\_tmax\_c} & raw & Primary heat driver of fire risk \\
\texttt{temp\_range} & $T_{\max} - T_{\min}$ & Diurnal range; proxy for continental/dry conditions \\
\texttt{tot\_prcp\_mm} & raw & Monthly precipitation; moisture suppressant \\
\texttt{dryness\_3m\_avg} & 3-month rolling mean & Drought accumulation; captures hydroclimate whiplash \\
\texttt{latitude} & geocoded & North--South geographic fire zone \\
\texttt{longitude} & geocoded & Coastal--inland gradient \\
\bottomrule
\end{tabular}
\label{tab:features}
\end{table}

The 3-month rolling dryness feature (\texttt{dryness\_3m\_avg}) is the most
important engineered feature. A single dry month presents limited fire risk;
the compounding effect of multiple consecutive dry and hot months --- described
in the challenge statement as ``hydroclimate whiplash'' --- is the physical
mechanism that \texttt{dryness\_3m\_avg} is designed to capture.

\subsection{Target Variable}

The target variable is the expected number of fires per ZIP-month:
\begin{equation}
    y = \log(1 + n_{\text{fires}})
    \label{eq:target}
\end{equation}
The log-transform stabilizes the heavily right-skewed distribution and
bounds the training signal. Predictions are inverted at inference time via
$\hat{n} = e^{\hat{y}} - 1$, clipped to $[0, \infty)$.

\subsection{Normalization}

All features are min-max normalized to $[0, 1]$ using statistics computed
exclusively on the training split (2018--2020) to prevent data leakage:
\begin{equation}
    x_{\text{norm}} = \frac{x - x_{\min}^{\text{train}}}{x_{\max}^{\text{train}} - x_{\min}^{\text{train}}}
\end{equation}
Values at inference time are clipped to $[0, 1]$. This normalization is
required for angle encoding in the quantum circuit, where each feature
is mapped to a rotation angle $\theta = \pi x_{\text{norm}} \in [0, \pi]$.
Scaler statistics are stored in \texttt{feature\_scaler\_stats.csv}.

% ══════════════════════════════════════════════════════════════════════════════
\section{Task 1A: Quantum Machine Learning Model}
\label{sec:qml}

\subsection{Motivation}
The challenge statement notes that ``the high degree of interconnectedness of
wildfire risk variables complicates relationships between features, which shows
promise for improvement through applying QML.'' Climate variables such as
temperature, precipitation, and dryness interact non-linearly and exhibit
geographic dependencies that may benefit from the expressive power of
quantum feature maps in high-dimensional Hilbert space.

\subsection{Circuit Architecture}

We implement a \textbf{Variational Quantum Circuit (VQC)} with 9 qubits
using PennyLane. The circuit consists of $L$ re-uploading layers, each
composed of three components:

\paragraph{Angle encoding.}
Each of the 9 normalized features is encoded as a single-qubit $R_y$ rotation:
\begin{equation}
    R_y(\pi x_i) = \begin{pmatrix} \cos(\pi x_i / 2) & -\sin(\pi x_i / 2) \\
    \sin(\pi x_i / 2) & \cos(\pi x_i / 2) \end{pmatrix}
\end{equation}
This maps each feature to a point on the Bloch sphere equator, placing all
inputs in the same rotation manifold.

\paragraph{Entangling layer.}
A ring of CNOT gates creates pairwise entanglement between adjacent qubits,
enabling the circuit to learn correlations between features (e.g.,
the joint signal from temperature and dryness):
\[
    \text{CNOT}(q_0, q_1),\ \text{CNOT}(q_1, q_2),\ \ldots,\ \text{CNOT}(q_8, q_0)
\]

\paragraph{Trainable rotation layer.}
Parameterized $R_y$ and $R_z$ rotations on each qubit provide the
trainable degrees of freedom:
\[
    R_z(\phi_i^{(l)}) R_y(\theta_i^{(l)}),\quad i = 0, \ldots, 8,\quad l = 1, \ldots, L
\]

\paragraph{Data re-uploading.}
The encoding-entangling-rotation sequence is repeated $L$ times, with the
same input data re-injected at each layer but with fresh trainable parameters.
This is analogous to depth in a classical neural network: a single encoding
layer can only represent linear functions of the features, while $L$
re-uploading layers enable the circuit to represent non-linear decision
boundaries~\cite{perez2020reup}.

\paragraph{Measurement.}
The expectation value $\langle Z_0 \rangle \in [-1, 1]$ is measured on qubit~0
and mapped to a non-negative prediction via a trainable affine output layer.

\subsection{Training Configuration}

\begin{table}[H]
\centering
\caption{QML model hyperparameters.}
\begin{tabular}{ll}
\toprule
\textbf{Parameter} & \textbf{Value} \\
\midrule
Qubits & 9 \\
Re-uploading layers $L$ & [TO BE FILLED] \\
Optimizer & Adam \\
Learning rate & [TO BE FILLED] \\
Batch size & [TO BE FILLED] \\
Epochs & [TO BE FILLED] \\
Loss function & Weighted MSE (fire rows: $10\times$ weight) \\
Backend & PennyLane \texttt{default.qubit} simulator \\
\bottomrule
\end{tabular}
\label{tab:qml_config}
\end{table}

\subsection{Resource Requirements}

\begin{table}[H]
\centering
\caption{Quantum resource estimates.}
\begin{tabular}{ll}
\toprule
\textbf{Resource} & \textbf{Value} \\
\midrule
Number of qubits & 9 \\
Circuit depth (per layer) & $\sim$27 gates ($9 R_y$ + $9$ CNOT + $9 R_y R_z$) \\
Total trainable parameters & $18 \times L$ \\
Shots (if sampling) & [TO BE FILLED] \\
Estimated wall-clock time (simulator) & [TO BE FILLED] \\
\bottomrule
\end{tabular}
\label{tab:resources}
\end{table}

\subsection{QML Results}
\label{sec:qml_results}

[TO BE FILLED after QML training is complete. Include: val/test MAE, RMSE,
Precision, Recall, F1, AUC-ROC using the same metrics as the classical
benchmark in Section~\ref{sec:classical}.]

% ══════════════════════════════════════════════════════════════════════════════
\section{Task 1B: Classical Benchmark and Comparison}
\label{sec:classical}

\subsection{Experimental Setup}

\paragraph{Temporal split.}
To prevent data leakage, a strict chronological split was applied:

\begin{table}[H]
\centering
\caption{Train / validate / test split.}
\begin{tabular}{lrrrr}
\toprule
\textbf{Split} & \textbf{Years} & \textbf{Rows} & \textbf{Fire rows} & \textbf{Fire rate} \\
\midrule
Train      & 2018--2020 & 93,348 & 1,026 & 1.10\% \\
Validate   & 2021       & 31,116 & 336   & 1.08\% \\
Test       & 2022       & 31,116 & 256   & 0.82\% \\
\midrule
\textbf{Total} & 2018--2022 & \textbf{155,580} & \textbf{1,618} & \textbf{1.04\%} \\
\bottomrule
\end{tabular}
\label{tab:split}
\end{table}

\paragraph{Sample weighting.}
To address the 99:1 class imbalance (Section~\ref{sec:imbalance}),
fire observations ($n_{\text{fires}} > 0$) were assigned a training sample
weight of 10 relative to no-fire observations.

\paragraph{Models.}
Five classical models were evaluated on identical features and splits:
(1)~Dummy Regressor (always predict zero, i.e., the na\"{i}ve baseline);
(2)~Ridge Regression; (3)~Random Forest (300 trees);
(4)~Gradient Boosting (300 estimators); (5)~XGBoost (300 estimators).

\paragraph{Metrics.}
All models are evaluated on the $\log(1 + n)$-scale predictions converted
back to the original \texttt{num\_fires} scale. Binary fire detection is
assessed at threshold $\hat{n} > 0.5$. Metrics reported:
Mean Absolute Error (MAE), Root Mean Squared Error (RMSE),
$R^2$, Precision, Recall, F1, and AUC-ROC.

\subsection{Validation Results}

\begin{table}[H]
\centering
\caption{Validation set (2021) results for all classical models.
         Best values per column in \textbf{bold}.}
\begin{tabular}{lrrrrrrr}
\toprule
\textbf{Model} & \textbf{MAE} & \textbf{RMSE} & \textbf{$R^2$} &
\textbf{Prec.} & \textbf{Recall} & \textbf{F1} & \textbf{AUC-ROC} \\
\midrule
Dummy (predict 0)  & \textbf{0.0124} & \textbf{0.1288} & $-$0.009 & 0.000 & 0.000 & 0.000 & 0.500 \\
Ridge Regression   & 0.0949 & 0.1584 & $-$0.527 & 0.000 & 0.000 & 0.000 & 0.728 \\
Random Forest      & 0.0511 & 0.1439 & $-$0.261 & 0.161 & 0.092 & 0.117 & 0.829 \\
XGBoost            & 0.0684 & 0.1569 & $-$0.499 & 0.177 & 0.158 & 0.167 & 0.820 \\
Gradient Boosting  & 0.0694 & 0.1549 & $-$0.460 & 0.163 & 0.131 & 0.145 & \textbf{0.837} \\
\bottomrule
\end{tabular}
\label{tab:val_results}
\end{table}

\subsection{Best Model: Test Set Results}

Gradient Boosting achieved the highest validation AUC-ROC and was
evaluated on the held-out 2022 test set (Table~\ref{tab:test_results}).

\begin{table}[H]
\centering
\caption{Gradient Boosting performance on the 2022 test set.}
\begin{tabular}{lr}
\toprule
\textbf{Metric} & \textbf{Value} \\
\midrule
MAE          & 0.0645 \\
RMSE         & 0.1426 \\
$R^2$        & $-$0.527 \\
Precision    & 0.242 \\
Recall       & 0.258 \\
F1           & 0.250 \\
\textbf{AUC-ROC} & \textbf{0.873} \\
Fire sensitivity     & 25.8\% \\
No-fire specificity  & 99.3\% \\
\bottomrule
\end{tabular}
\label{tab:test_results}
\end{table}

\subsection{Interpretation of Metrics}

\paragraph{Why $R^2 < 0$?}
A negative $R^2$ does not indicate a model worse than a random predictor.
For zero-inflated data, $R^2$ compares the model to the global mean
prediction; since the global mean is non-zero but most true values are
exactly zero, any model that correctly predicts mostly zeros will have
low or negative $R^2$. AUC-ROC is the appropriate primary metric here
because it measures \emph{ranking quality}: whether fire ZIP-months
are scored higher than no-fire ZIP-months regardless of the absolute
prediction scale.

\paragraph{AUC-ROC of 0.873.}
This means that if one fire ZIP-month and one no-fire ZIP-month are drawn
at random, the Gradient Boosting model ranks the fire ZIP-month higher
with 87.3\% probability. This is a strong discriminative signal for
insurance pricing applications.

\paragraph{Low F1 is structural.}
With a 1\% base rate, a threshold of $\hat{n} > 0.5$ yields low F1 regardless
of model quality. For insurance applications, the model output is used as
a continuous risk score (not a hard binary), so F1 is less relevant than
AUC-ROC and the calibration of the continuous score.

\subsection{2023 Risk Score Generation}

Monthly fire predictions for all 2,593 ZIP codes across all 12 months of
2023 were generated using historical ZIP$\times$month averages as synthetic
2023 features (since 2023 climate observations are not available in the
provided dataset). The annual risk score for each ZIP code is:
\begin{equation}
    \text{RiskScore}_{z} = \sum_{m=1}^{12} \hat{n}_{z,m}
    \label{eq:risk}
\end{equation}
Table~\ref{tab:top10} shows the ten highest-risk ZIP codes for 2023.

\begin{table}[H]
\centering
\caption{Top 10 highest-risk ZIP codes for 2023 (Gradient Boosting model).}
\begin{tabular}{rrl}
\toprule
\textbf{ZIP} & \textbf{Annual Risk Score} & \textbf{Region} \\
\midrule
93453 & 13.52 & Santa Maria / SLO County \\
93461 & 13.20 & Paso Robles area \\
93446 & 9.30  & Paso Robles \\
93210 & 8.99  & Coalinga / Fresno County \\
96130 & 7.31  & Lassen County (Northern CA) \\
93001 & 6.71  & Ventura County \\
93426 & 6.60  & Monterey County \\
93280 & 6.31  & Wasco / Kern County \\
96064 & 6.28  & Siskiyou County (Northern CA) \\
93432 & 6.13  & Arroyo Grande area \\
\bottomrule
\end{tabular}
\label{tab:top10}
\end{table}

The predicted high-risk ZIP codes align geographically with California's
historically fire-prone regions: the Central Coast (San Luis Obispo and
Santa Barbara counties), the Central Valley interface, and Northern California
mountain counties, providing face validity for the model.

\subsection{QML vs.\ Classical Comparison}
\label{sec:comparison}

[TO BE FILLED after QML model is trained. Structure: side-by-side table of
all metrics for best classical (Gradient Boosting) vs.\ QML model on the
2022 test set. Discuss: (1) whether QML improves AUC-ROC or recall;
(2) trade-off between performance and circuit depth/training cost;
(3) advantages of the quantum feature map for capturing non-linear
feature interactions.]

\paragraph{Advantages of the QML approach (anticipated).}
\begin{itemize}[leftmargin=*]
    \item The quantum feature map embeds 9-dimensional input into a
          $2^9 = 512$-dimensional Hilbert space, enabling non-linear
          decision boundaries that are exponentially costly to reproduce
          classically.
    \item Entangling gates between climate feature pairs (e.g.,
          temperature--dryness, lat--lon) explicitly model the
          interdependencies described in the challenge statement.
    \item As the number of relevant features grows with climate complexity,
          QML circuits scale more efficiently in expressivity per parameter
          than classical polynomial feature expansions.
\end{itemize}

\paragraph{Disadvantages of the QML approach.}
\begin{itemize}[leftmargin=*]
    \item Training is significantly slower than classical models on current
          simulators (no hardware acceleration for gradient computation).
    \item Barren plateaus in the parameter landscape can impede convergence
          for deep circuits.
    \item Current quantum hardware introduces gate noise that degrades
          prediction quality; results on real devices may underperform
          the noise-free simulator.
    \item The 1\% fire rate makes gradient signals sparse during training,
          which is exacerbated by the shot noise inherent to quantum
          measurement.
\end{itemize}

% ══════════════════════════════════════════════════════════════════════════════
\section{Task 2: Insurance Premium Modeling}
\label{sec:task2}

[TO BE FILLED. The annual risk scores from Equation~\eqref{eq:risk} feed
directly into a time series model predicting 2021 premiums from 2018--2020
data. Use \texttt{annual\_risk\_scores\_2023.csv} as the fire risk input
column alongside the insurance dataset.]

% ══════════════════════════════════════════════════════════════════════════════
\section{Conclusions}
\label{sec:conclusions}

We have developed a complete data pipeline, feature engineering framework,
and classical benchmark for California wildfire risk prediction.
The cleaned dataset covers 2,593 ZIP codes across 2018--2022 with no
missing values in any modeled feature. The Gradient Boosting classical model
achieves an AUC-ROC of 0.873 on the 2022 held-out test set, providing a
strong benchmark for the QML model to be compared against.

The 9-qubit VQC with angle encoding and data re-uploading is designed to
operate within the circuit budget of current simulators (AWS Braket,
IBM Quantum) while providing the expressivity needed to capture the
non-linear interactions between climate, geography, and historical fire
patterns. The annual risk scores produced by both models are structured
for direct ingestion by the Task~2 insurance premium model.

% ── References ────────────────────────────────────────────────────────────────
\begin{thebibliography}{9}

\bibitem{perez2020reup}
A.~P\'{e}rez-Salinas, A.~Cervera-Lierta, E.~Gil-Fuster, and J.~I.~Latorre,
``Data re-uploading for a universal quantum classifier,''
\textit{Quantum}, vol.~4, p.~226, 2020.
\href{https://doi.org/10.22331/q-2020-02-06-226}{doi:10.22331/q-2020-02-06-226}

\bibitem{cerezo2021vqe}
M.~Cerezo et al.,
``Variational quantum algorithms,''
\textit{Nature Reviews Physics}, vol.~3, pp.~625--644, 2021.

\bibitem{schuld2021supervised}
M.~Schuld and N.~Killoran,
``Quantum machine learning in feature Hilbert spaces,''
\textit{Physical Review Letters}, vol.~122, no.~4, p.~040504, 2019.

\bibitem{wri2025fires}
J.~MacCarthy et al.,
``The Latest Data Confirms: Forest Fires Are Getting Worse,''
World Resources Institute, July 2025.
\url{https://www.wri.org/insights/global-trends-forest-fires}

\end{thebibliography}

\end{document}
